\section{Wyznaczanie poziomu mocy akustycznej}
\subsection{Rozmieszczenie punktów i przebieg pomiarów}
Wykonano pięć pomiarów przy pracującym źródle na środku pomieszczenia
oraz po trzy przy źródle w rogu i przy ścianie. Rejestrację P38 przypisano
do tła; nie wchodzi ona do średniej poziomów pracującego źródła.
Określenia ścian w tabeli zachowują nazwy z notatek i wskazują kierunki
punktów w ustawieniu środkowym, a nie osobne położenia odkurzacza.
Dla środka zidentyfikowano punkt nad źródłem i cztery kierunki boczne.
Dla rogu i ściany zachowano kolejność rejestracji; dokładnych współrzędnych
mikrofonów nie odtworzono. Nie dopisywano punktów na podstawie podobieństwa
poziomu dźwięku.

\subsection{Zestawienie pomiarów}
Wartości $L_{Aeq}$ odczytano z plików SVL miernika SVAN 971.
Każda z przypisanych rejestracji trwała \qty{15}{\second}; zapis sekundowy
obejmuje 15 rekordów. Profile mają kolejno korekcje A, C i Z.
Sprawdzono zgodność średnich energetycznych rekordów z wynikami zbiorczymi
każdego profilu, z tolerancją \qty{0.011}{\decibel}, wynikającą z rozdzielczości zapisu.
Interpretacja formatu opiera się na instrukcji producenta \autocite{svan971}.

Wartości z notatek odpowiadają po zaokrągleniu poziomowi ekspozycji $L_{AE}$,
a nie poziomowi równoważnemu $L_{Aeq}$:
\begin{equation}
L_{AE}=L_{Aeq}+10\lg(T/\qty{1}{\second}).
\end{equation}
Dla $T=\qty{15}{\second}$ różnica wynosi około \qty{11.76}{\decibel}.
Przykładowo P37 zawiera $L_{Aeq}=\qty{82.70}{\decibel}$ oraz
$L_{AE}=\qty{94.46}{\decibel}$. W obliczeniach wykorzystano $L_{Aeq}$.

\begin{table}[H]
\centering\small
\caption{Poziomy równoważne A z SVL; P38 stanowi tło.}
\begin{tabularx}{\linewidth}{@{}l l L r@{}}
\toprule
Położenie & Plik & Opis punktu & $L_{Aeq}$ (dB)\\
\midrule
Środek & P37 & nad odkurzaczem & \num{82.7} \\
Środek & P38 & tło & \num{27.65} \\
Środek & P39 & ściana paleta & \num{84.34} \\
Środek & P40 & ściana drzwi & \num{84.13} \\
Środek & P41 & ściana metal & \num{82.27} \\
Środek & P42 & ściana apteczka & \num{83.88} \\
Róg & P50 & do ustalenia & \num{85.44} \\
Róg & P51 & do ustalenia & \num{84.77} \\
Róg & P52 & do ustalenia & \num{87.03} \\
Ściana & P53 & do ustalenia & \num{85.01} \\
Ściana & P54 & do ustalenia & \num{84.0} \\
Ściana & P55 & do ustalenia & \num{84.23} \\
\bottomrule

\end{tabularx}
\end{table}

\subsection{Uśrednianie i korekcje}
Dla jednakowych wag punktów średnią energetyczną oblicza się jako:
\begin{equation}
\overline L=10\lg\left(\frac{1}{n}\sum_{i=1}^n10^{0{,}1L_i}\right).
\end{equation}
\begin{table}[H]
\centering
\caption{Średnie energetyczne $L_{Aeq}$ z SVL, przy jednakowych wagach.}
\begin{tabular}{@{}l c r@{}}
\toprule
Położenie & Liczba punktów źródła & $\overline L_A$ (dB)\\
\midrule
Środek & 5 & \num{83.54} \\
Róg & 3 & \num{85.85} \\
Ściana & 3 & \num{84.44} \\
\bottomrule

\end{tabular}
\end{table}
Są to średnie poziomów ciśnienia, przed korekcjami i przeliczeniem na moc.
Jeśli punkty reprezentują różne pola, należy zastąpić jednakowe wagi
udziałami powierzchni przypisanych punktom.

Dla poziomów A stosuje się
$\Delta L_A=\overline L_{A,\text{źr+tło}}-\overline L_{A,\text{tło}}$
oraz $K_{1A}=-10\lg(1-10^{-0{,}1\Delta L_A})$ w zakresie wskazanym
w konspekcie. Dla $\Delta L_A>\qty{10}{\decibel}$ konspekt przyjmuje
$K_{1A}=0$. Przy różnicy poniżej \qty{3}{\decibel} kryterium tła nie jest
spełnione. Poziom $L_{Aeq}$ tła P38 wynosi \qty{27.65}{\decibel} i jest o ponad \qty{50}{\decibel} niższy od średniej
dla środka, lecz pojedynczy pomiar tła nie potwierdza jego rozkładu na
całej powierzchni ani warunków przy ścianie i w rogu.

Seria przy ścianie zawiera trzy potwierdzone punkty. Nie zidentyfikowano
czwartego punktu na podstawie zachowanych plików. Jest to ograniczenie
próbkowania powierzchni, którego nie usuwa samo uśrednienie wyników.

\subsection{Przeliczenie na poziom mocy}
W przyjętym modelu obliczeniowym zastosowano zależność:
\begin{equation}
L_{WA}=\overline L_{A,\text{źr+tło}}-K_{1A}-K_{2A}
+10\lg\left(\frac{S}{S_0}\right),\qquad S_0=\qty{1}{\metre\squared}.
\end{equation}

Wyniki zestawiono przy wspólnym tle P38, jednakowych wagach punktów
oraz geometrii opisanej w części 1. Przeniesienie pomiaru tła na pozostałe
ustawienia jest założeniem opracowania; duża różnica poziomów wskazuje
na pomijalny wpływ tego tła przy założeniu jego stałości.
\begin{table}[H]
\centering\small
\caption{Oszacowanie poziomu mocy A; wynik przy ścianie oparty na trzech punktach.}
\begin{tabular}{@{}l c r r r r r@{}}
\toprule
Położenie & $n$ & $S$ (m$^2$) & $\overline L_A$ & $K_{1A}$ & $K_{2A}$ & $L_{WA}$\\
 & & & \multicolumn{4}{c}{dB}\\
\midrule
Środek & 5 & \num{16.79} & \num{83.54} & \num{0.00} & \num{5.68} & \num{90.11} \\
Ściana & 3 & \num{9.02} & \num{84.44} & \num{0.00} & \num{3.89} & \num{90.10} \\
Róg & 3 & \num{5.10} & \num{85.85} & \num{0.00} & \num{2.60} & \num{90.33} \\
\bottomrule

\end{tabular}
\end{table}
Dla środka przykładowe przeliczenie daje:
\begin{equation}
L_{WA}=83{,}54-0-5{,}68+10\lg(16{,}7912)
\approx\qty{90.11}{\decibel}.
\end{equation}
Różnica największego i najmniejszego oszacowania wynosi około
\qty{0.23}{\decibel}. Nie stanowi to dowodu tak małej niepewności pomiaru:
wszystkie wyniki korzystają z tego samego przybliżenia chłonności i tła.
